% Pursuing Stacks (cotes 134-2 to 134-8): archive cote 134-5, site page 61 (PDF page 62), typescript page 248.
% typescript. Draft compared: ps4.tex:3076-3140.
% AI-transcribed by Claude Opus 5.5 (claude-opus-5-5, adaptive thinking, effort high), Message Batches API, batch ps2 (msgbatch_01NVwr1Yuo8qJf5rykeGBXPK), 2026-11-12; not yet checked by eye.
\scanpage{134-5/61}
\mspage{\deleted{248} \hand{248}}

\xout{Max} Let now
\inserted[typed]{$(M,M_c)$} \inserted[typed]{be a}\xout{xxxxxxxxxxx2xxLetx(M,Mx)xxxxxx}contractibility structure, such that
there exist a \uline{small} \inserted{full} subcategory $C$ of $M$ \xout{whosexobjectsxarexcontractible}
\over{W}{w}hich generates the contractibility structure, i.e. such that \xout{xx}
(a) the objects of $C$ are ``contractible'' i.e. are in $M_c$ and (b) any
object of $M_c$ \deleted{which} is $C$-contractible\deleted{\unsure{s}} (i.e. contractible for the homotopy
interval structure admitting the intervals made up with objects of $C$ as
a generating family)\hand{,} \deleted{is in $M_c$} . Independently of any smallness assump-%
tion upon $C$, we gave \unsure{in} section 51 D) \xout{the} (p.118), under the name
of ``basic assumption'' (Bas 4), the ns condition on a full subcategory
$C$ of a given category $M$, in order that $C$ can be viewed as a generating
set of contractible objects, for a suitable contractibility structure
\xout{Mx XX}
\[ M_c \subset \mathrm{Ob} M \]
in $M$ (which is uniquely defined by $C$). It turns out that in case $C$ \xout{co}con-%
\xout{stable} tains a final object of $M$, and is stable \inserted{under} \deleted{by} binary products in
$M$, the condition (Bas 4) \xout{doesxnot} depends only \inserted{upon} \deleted{on} the category structure
of $C$ , and not upon the particular way this category is embedded in
an\over{o}{o}ther one $M$.

In loc. cit. we did not impose, when defining a contractibility
structure, a condition that the\over{n}{r} should exist a small set of generators
for the structure. From now on, we'll assume that the (possibly large)
categories $M$ we are working with are $\underline{U}$-categories, and that ``contracti-%
bility structure'' means ``$\underline{U}$-contractibility structure'', namely one \inserted{such that} \deleted{admit-}
\inserted{M admit} \deleted{ting} a small full subcategory $C$, generating the structure.

We recall too that in the definition of a contractibility structure
$(M,M_c)$, it has always been understood that $M$ is stable under finite pro-%
ducts.

\uline{Proposition} 2. With the conventions above, let $(M,M_c)$ be any contrac-%
tibility structure. Then:

a) The full subcategory $M_c$ of $M$ generates an asphericity structure
$M_{as}$ in $M$ (cf. prop.1).

b) \uline{Any small full} subcategory $C$ of $M$ which generates the contractibi
lity structure\inserted{and such that $C^{\wedge}$ be totally aspheric,} \uline{generates} the asphericity structure $M_{as}$\hand{,} i.e. the inclu-%
sion functor $i: C \to M$ is $M_{as}$-aspheric.

c) The asphericity structure $M_{as}$ is totally aspheric (cf. \over{P}{p}rop.5
p. 239).

\margin{61}
